\documentclass[10pt,a4paper]{article}

\usepackage{geometry}
\geometry{
  left=15mm,
  right=15mm,
  top=15mm,
  bottom=15mm
}
\usepackage{amsmath}
\usepackage{amssymb}
\usepackage{enumitem}

\setlist[enumerate,1]{label=(\arabic*),leftmargin=2.2em}
\setlist[enumerate,2]{label=(\alph*),leftmargin=2.2em}

\begin{document}

\begin{center}
  {\Large \textbf{The University of Tokyo 2024}}\\[2mm]
  {\large \textbf{Problems 1--3}}
\end{center}

\section*{Problem 1}

Let $\mathbb{R}^3$ be the set of the three-dimensional real column vectors and $\mathbb{R}^{3\times 3}$ be the set of the three-by-three real matrices. Let $n_1,n_2,n_3\in\mathbb{R}^3$ be linearly independent unit-length vectors and $n_4\in\mathbb{R}^3$ be a unit-length vector not parallel to $n_1,n_2,$ or $n_3$. Let $A$ and $B$ be square matrices defined as
\[
  A=
  \begin{pmatrix}
    n_1^T-n_2^T\\
    n_1^T-n_3^T\\
    n_3^T-n_4^T
  \end{pmatrix},
  \qquad
  B=\sum_{i=1}^4 n_in_i^T.
\]
Here, $X^T$ and $x^T$ denote the transpose of a matrix $X$ and a vector $x$, respectively. Answer the following questions.

\begin{enumerate}
  \item Find the condition for $n_4$ such that the rank of $A$ is three.

  \item In the three-dimensional Euclidean space $\mathbb{R}^3$, consider four planes
        \[
          \Pi_i=\{x\in\mathbb{R}^3\mid n_i^Tx-d_i=0\}
        \]
        where $d_i$ is a real number, and $i=1,2,3,4$, that satisfy the following three conditions: (i) the rank of $A$ is three, (ii)
        \[
          \Omega=\{x\in\mathbb{R}^3\mid n_i^Tx-d_i\geq 0,\ i=1,2,3,4\}
        \]
        is not the empty set, and (iii) there exists a sphere C $(\subset\Omega)$ to which $\Pi_i$ $(i=1,2,3,4)$ are tangent. The position vector of the center of C is represented by $A^{-1}u$ using a vector $u\in\mathbb{R}^3$. Express $u$ using $d_i$ $(i=1,2,3,4)$.

  \item Show that $B$ is a positive definite symmetric matrix.

  \item Consider the point P from which the sum of squared distances to four planes
        \[
          \{x\in\mathbb{R}^3\mid n_i^Tx-d_i=0\}
        \]
        where $d_i$ is a real number, and $i=1,2,3,4$, is minimized. The position vector of P is represented by $B^{-1}v$ using a vector $v\in\mathbb{R}^3$. Express $v$ using $n_i$ and $d_i$ $(i=1,2,3,4)$.

  \item Let $l_i$ be a straight line through a point $Q_i$, the position vector of which is $x_i\in\mathbb{R}^3$, parallel to $n_i$ $(i=1,2,3)$ in $\mathbb{R}^3$. Let $R_i$ be the orthogonal projection of an arbitrary point $R$, the position vector of which is $y\in\mathbb{R}^3$, onto $l_i$. The position vector of $R_i$ is represented by
        \[
          y-W_i(y-x_i)
        \]
        using a matrix $W_i\in\mathbb{R}^{3\times 3}$. The identity matrix is denoted by $I\in\mathbb{R}^{3\times 3}$.
        \begin{enumerate}
          \item Express $W_i$ using $n_i$ and $I$.
          \item Show that $W_i^TW_i=W_i$.
          \item Consider a plane
                \[
                  \Sigma=\{x\in\mathbb{R}^3\mid a^Tx=b\}
                \]
                where $a\in\mathbb{R}^3$ is a non-zero vector, and $b$ is a real number. Let $S\in\Sigma$ be the point from which the sum of squared distances to $l_1,l_2,$ and $l_3$ is minimized. When $n_1,n_2,$ and $n_3$ are orthogonal to each other, the position vector of $S$ is represented by
                \[
                  \left(I-\frac{aa^T}{a^Ta}\right)w+\frac{ab}{a^Ta}
                \]
                using a vector $w\in\mathbb{R}^3$ which is independent of $a$ and $b$. Express $w$ using $W_i$ and $x_i$ $(i=1,2,3)$.
        \end{enumerate}
\end{enumerate}

\section*{Problem 2}

Consider a function $f(s)$ defined by the following integral for positive real numbers $s$:
\[
  f(s)=\int_0^\infty t^{s-1}\exp(-t)\,dt.
\]
Answer the following questions. You may answer without showing that the above integral converges.

\begin{enumerate}
  \item Find the value of $f(1)$.

  \item The inequality $\exp(t)>t^n/n!$ holds for any positive real number $t$ and non-negative integer $n$.
        \begin{enumerate}
          \item For positive real numbers $s$, show the following inequality:
                \[
                  \int_0^1 t^{s-1}\exp(-t)\,dt<\frac{1}{s}.
                \]
          \item When $n>s>0$, show that the following inequality holds for any real number $c$ that satisfies $c>1$:
                \[
                  \int_1^c t^{s-1}\exp(-t)\,dt<\frac{n!}{n-s}.
                \]
        \end{enumerate}

  \item When the second-order derivative of $f(s)$ is expressed as
        \[
          \frac{d^2f(s)}{ds^2}=\int_0^\infty g(t,s)\exp(-t)\,dt,
        \]
        find a function $g(t,s)$. You may answer without showing that the order of differentiation and integration can be exchanged.

  \item Find the value of $D$ defined as
        \[
          D=\int_0^\infty(\log t)^2\exp(-t)\,dt
          -
          \left\{\int_0^\infty(\log t)\exp(-t)\,dt\right\}^2.
        \]
        Here, you may use the fact that the following relation holds:
        \[
          \left.\frac{d^2}{ds^2}\log f(s)\right|_{s=1}=\frac{\pi^2}{6}.
        \]

  \item Define a function $p(r)$ for positive real numbers $r$ and $\alpha$ as
        \[
          p(r)=\frac{r}{\alpha}\exp\left(-\frac{r^2}{2\alpha}\right).
        \]
        Find the value of $S$ defined as
        \[
          S=\int_0^\infty(\log r)^2p(r)\,dr
          -
          \left\{\int_0^\infty(\log r)p(r)\,dr\right\}^2.
        \]
\end{enumerate}

\section*{Problem 3}

Consider a particle moving on the coordinate plane, and denote the location of the particle at time $t\in\{0,1,2,\ldots\}$ by $(X_t,Y_t)$. The initial location of the particle is $(X_0,Y_0)=(0,0)$. Also, if $(X_t,Y_t)=(a,b)$, then $(X_{t+1},Y_{t+1})=(a+1,b)$ with probability $p$, $(X_{t+1},Y_{t+1})=(a,b+1)$ with probability $q$, and $(X_{t+1},Y_{t+1})=(a,b)$ with probability $1-p-q$. Here, it is assumed that $p,q>0$, $p+q<1$, and the movements of the particle at different time points are independent. Let $(X,Y)$ denote the location of the particle such that
\[
  (X_{t+1},Y_{t+1})=(X_t,Y_t)
\]
for the first time. Answer the following questions.

\begin{enumerate}
  \item Show that the probability that $(X,Y)=(1,2)$ is $3pq^2(1-p-q)$.

  \item For non-negative integers $n$, find the probability that $X+Y=n$.

  \item For non-negative integers $n$, let $f_n$ denote the probability that $X=n$.
        \begin{enumerate}
          \item Find $f_0$.
          \item Express the probability that $X\geq n+1$ given the condition $X\geq n$, using $f_0$.
          \item Show that $f_n=(1-f_0)^nf_0$.
        \end{enumerate}

  \item Express the expectation of $X$ using $p$ and $q$.

  \item Express the correlation coefficient between $X$ and $Y$
        \[
          \frac{E[(X-\mu_X)(Y-\mu_Y)]}
          {\sqrt{E[(X-\mu_X)^2]E[(Y-\mu_Y)^2]}}
        \]
        using $p$ and $q$, where $\mu_X=E[X]$ denotes the expectation of $X$ and $\mu_Y=E[Y]$ denotes the expectation of $Y$.
\end{enumerate}

\newpage

\begin{center}
  {\Large \textbf{Solutions 1--3}}
\end{center}

\section*{Solution 1}

\begin{enumerate}
  \item Since $n_1,n_2,n_3$ are linearly independent unit vectors, they form a basis of $\mathbb{R}^3$. Write
        \[
          n_4=an_1+bn_2+cn_3.
        \]
        With respect to this basis, the three rows of $A$ correspond to
        \[
          (1,-1,0),\qquad (1,0,-1),\qquad (-a,-b,1-c).
        \]
        Their determinant is
        \[
          \begin{vmatrix}
            1&-1&0\\
            1&0&-1\\
            -a&-b&1-c
          \end{vmatrix}
          =1-a-b-c.
        \]
        Therefore
        \[
          \operatorname{rank}A=3
          \iff
          1-a-b-c\neq 0.
        \]
        Equivalently,
        \[
          \boxed{
          n_4^T(n_1+n_2+n_3)\neq 1.}
        \]

  \item Let $z$ be the center of the sphere and let its radius be $r$. Since the sphere is tangent to all four planes and lies in $\Omega$,
        \[
          n_i^Tz-d_i=r
          \qquad (i=1,2,3,4).
        \]
        Subtracting these equations gives
        \[
          (n_1-n_2)^Tz=d_1-d_2,
        \]
        \[
          (n_1-n_3)^Tz=d_1-d_3,
        \]
        and
        \[
          (n_3-n_4)^Tz=d_3-d_4.
        \]
        Hence $Az=u$, where
        \[
          \boxed{
          u=
          \begin{bmatrix}
            d_1-d_2\\
            d_1-d_3\\
            d_3-d_4
          \end{bmatrix}.}
        \]
        Since $\operatorname{rank}A=3$, the center is $z=A^{-1}u$.

  \item The matrix $B$ is symmetric because
        \[
          (n_in_i^T)^T=n_in_i^T
        \]
        for every $i$. For any nonzero $x\in\mathbb{R}^3$,
        \[
          x^TBx
          =
          \sum_{i=1}^4 x^Tn_in_i^Tx
          =
          \sum_{i=1}^4(n_i^Tx)^2.
        \]
        Since $n_1,n_2,n_3$ are linearly independent, no nonzero $x$ can be orthogonal to all of $n_1,n_2,n_3$. Thus
        \[
          x^TBx>0
        \]
        for every nonzero $x$. Therefore $B$ is positive definite symmetric.

  \item The squared distance from $x$ to the plane $n_i^Tx-d_i=0$ is
        \[
          (n_i^Tx-d_i)^2,
        \]
        because $n_i$ is a unit vector. Thus we minimize
        \[
          F(x)=\sum_{i=1}^4(n_i^Tx-d_i)^2.
        \]
        Differentiating gives
        \[
          \nabla F(x)=2\sum_{i=1}^4(n_i^Tx-d_i)n_i.
        \]
        Setting this equal to zero yields
        \[
          \left(\sum_{i=1}^4 n_in_i^T\right)x
          =
          \sum_{i=1}^4d_in_i.
        \]
        Therefore the minimizing point is $x=B^{-1}v$, where
        \[
          \boxed{
          v=\sum_{i=1}^4d_in_i.}
        \]

  \item
        \begin{enumerate}
          \item The orthogonal projection of $y$ onto the line through $x_i$ in the direction $n_i$ is
                \[
                  x_i+n_in_i^T(y-x_i).
                \]
                Since this is also written as
                \[
                  y-W_i(y-x_i),
                \]
                we obtain
                \[
                  \boxed{
                  W_i=I-n_in_i^T.}
                \]

          \item Since $n_i$ is a unit vector,
                \[
                  W_i^T=I-n_in_i^T=W_i
                \]
                and
                \[
                  W_i^2
                  =
                  (I-n_in_i^T)^2
                  =
                  I-2n_in_i^T+n_i(n_i^Tn_i)n_i^T
                  =
                  I-n_in_i^T
                  =
                  W_i.
                \]
                Hence
                \[
                  \boxed{W_i^TW_i=W_i.}
                \]

          \item The squared distance from $s$ to $l_i$ is
                \[
                  \|W_i(s-x_i)\|^2.
                \]
                By part (b), this is
                \[
                  (s-x_i)^TW_i(s-x_i).
                \]
                Therefore the sum of squared distances is minimized, without the plane constraint, when
                \[
                  \left(\sum_{i=1}^3W_i\right)s
                  =
                  \sum_{i=1}^3W_ix_i.
                \]
                If $n_1,n_2,n_3$ are orthogonal unit vectors, then
                \[
                  \sum_{i=1}^3W_i
                  =
                  \sum_{i=1}^3(I-n_in_i^T)
                  =
                  3I-I
                  =
                  2I.
                \]
                Hence the unconstrained minimizer is
                \[
                  w=\frac12\sum_{i=1}^3W_ix_i.
                \]
                The constrained minimizer on $a^Ts=b$ is the orthogonal projection of $w$ onto $\Sigma$, namely
                \[
                  \left(I-\frac{aa^T}{a^Ta}\right)w+\frac{ab}{a^Ta}.
                \]
                Thus
                \[
                  \boxed{
                  w=\frac12\sum_{i=1}^3W_ix_i.}
                \]
        \end{enumerate}
\end{enumerate}

\section*{Solution 2}

\begin{enumerate}
  \item Since
        \[
          f(1)=\int_0^\infty e^{-t}\,dt,
        \]
        we have
        \[
          \boxed{f(1)=1.}
        \]

  \item
        \begin{enumerate}
          \item Since $e^{-t}<1$ for $t>0$,
                \[
                  \int_0^1t^{s-1}e^{-t}\,dt
                  <
                  \int_0^1t^{s-1}\,dt
                  =
                  \frac1s.
                \]

          \item From
                \[
                  e^t>\frac{t^n}{n!},
                \]
                we get
                \[
                  e^{-t}<\frac{n!}{t^n}.
                \]
                Therefore
                \[
                  \int_1^c t^{s-1}e^{-t}\,dt
                  <
                  n!\int_1^c t^{s-n-1}\,dt.
                \]
                Since $n>s$,
                \[
                  n!\int_1^c t^{s-n-1}\,dt
                  =
                  \frac{n!}{n-s}\left(1-c^{s-n}\right)
                  <
                  \frac{n!}{n-s}.
                \]
        \end{enumerate}

  \item Differentiating twice with respect to $s$ gives
        \[
          \frac{d^2}{ds^2}t^{s-1}
          =
          (\log t)^2t^{s-1}.
        \]
        Hence
        \[
          \boxed{
          g(t,s)=(\log t)^2t^{s-1}.}
        \]

  \item We have
        \[
          f'(s)=\int_0^\infty(\log t)t^{s-1}e^{-t}\,dt,
        \]
        and
        \[
          f''(s)=\int_0^\infty(\log t)^2t^{s-1}e^{-t}\,dt.
        \]
        Since $f(1)=1$,
        \[
          \left.\frac{d^2}{ds^2}\log f(s)\right|_{s=1}
          =
          f''(1)-\{f'(1)\}^2
          =
          D.
        \]
        Using the given relation,
        \[
          \boxed{D=\frac{\pi^2}{6}.}
        \]

  \item Put
        \[
          u=\frac{r^2}{2\alpha}.
        \]
        Then
        \[
          du=\frac{r}{\alpha}\,dr,
        \]
        so $u$ has density $e^{-u}$ on $(0,\infty)$. Also,
        \[
          \log r=\frac12\log(2\alpha)+\frac12\log u.
        \]
        Adding a constant does not change variance, and multiplying by $1/2$ multiplies variance by $1/4$. Therefore, using the result of part (4),
        \[
          S=\frac14\cdot\frac{\pi^2}{6}.
        \]
        Hence
        \[
          \boxed{S=\frac{\pi^2}{24}.}
        \]
\end{enumerate}

\section*{Solution 3}

Let
\[
  r=1-p-q.
\]
The particle moves right with probability $p$, moves upward with probability $q$, and stops for the first time with probability $r$.

\begin{enumerate}
  \item To end at $(1,2)$ before the first stop, the particle must make exactly one right move and two upward moves in the first three steps, in any order, and then stop. Hence
        \[
          P((X,Y)=(1,2))
          =
          \binom31pq^2r
          =
          \boxed{3pq^2(1-p-q).}
        \]

  \item The event $X+Y=n$ means that the first $n$ steps are moves and the next step is the first stop. Therefore
        \[
          \boxed{
          P(X+Y=n)=(p+q)^n(1-p-q).}
        \]

  \item
        \begin{enumerate}
          \item The event $X=0$ means no right move occurs before the first stop. At each step, either an upward move occurs, with probability $q$, or the first stop occurs, with probability $r$. Thus
                \[
                  f_0=r(1+q+q^2+\cdots)
                  =
                  \frac{r}{1-q}.
                \]
                Hence
                \[
                  \boxed{
                  f_0=\frac{1-p-q}{1-q}.}
                \]

          \item The number $X$ has a geometric distribution:
                \[
                  f_n
                  =
                  \sum_{m=0}^\infty
                  \binom{n+m}{n}p^nq^mr
                  =
                  \frac{r}{1-q}
                  \left(\frac{p}{1-q}\right)^n.
                \]
                Therefore
                \[
                  P(X\geq n+1\mid X\geq n)
                  =
                  \frac{p}{1-q}.
                \]
                Since
                \[
                  1-f_0
                  =
                  1-\frac{1-p-q}{1-q}
                  =
                  \frac{p}{1-q},
                \]
                we get
                \[
                  \boxed{
                  P(X\geq n+1\mid X\geq n)=1-f_0.}
                \]

          \item From the result above,
                \[
                  \boxed{
                  f_n=(1-f_0)^nf_0.}
                \]
        \end{enumerate}

  \item Since $X$ has the geometric distribution
        \[
          f_n=(1-f_0)^nf_0,
        \]
        its expectation is
        \[
          E[X]=\frac{1-f_0}{f_0}.
        \]
        Substituting
        \[
          f_0=\frac{1-p-q}{1-q}
        \]
        gives
        \[
          \boxed{
          E[X]=\frac{p}{1-p-q}.}
        \]

  \item The joint probability generating function of $(X,Y)$ is
        \[
          G(s,t)=E[s^Xt^Y]
          =
          \sum_{n=0}^\infty r(ps+qt)^n
          =
          \frac{r}{1-ps-qt}.
        \]
        Hence
        \[
          E[X]=\frac{p}{r},
          \qquad
          E[Y]=\frac{q}{r},
        \]
        \[
          E[XY]
          =
          \left.\frac{\partial^2G}{\partial s\,\partial t}\right|_{s=t=1}
          =
          \frac{2pq}{r^2}.
        \]
        Therefore
        \[
          \operatorname{Cov}(X,Y)
          =
          E[XY]-E[X]E[Y]
          =
          \frac{pq}{r^2}.
        \]
        Similarly,
        \[
          \operatorname{Var}(X)=\frac{p(1-q)}{r^2},
          \qquad
          \operatorname{Var}(Y)=\frac{q(1-p)}{r^2}.
        \]
        Thus the correlation coefficient is
        \[
          \frac{\operatorname{Cov}(X,Y)}
          {\sqrt{\operatorname{Var}(X)\operatorname{Var}(Y)}}
          =
          \boxed{
          \sqrt{\frac{pq}{(1-p)(1-q)}}.}
        \]
\end{enumerate}

\end{document}
